\documentclass[fontsize=11pt]{article}

\usepackage[margin=1in]{geometry}
\usepackage{graphicx}
\usepackage{wrapfig}
\usepackage{subcaption}
\usepackage{amsmath} 
\usepackage{siunitx}
\usepackage{booktabs}
\usepackage{gensymb}
\usepackage{hyperref}
\usepackage[export]{adjustbox}
\usepackage{anyfontsize}

\begin{document}

	\graphicspath{{./figures/}}

	%BJC: check address, phone number and fax (?) number
	\begin{figure}
	\centering
	\begin{minipage}{0.37\textwidth}
	\includegraphics[width=2.14in,left]{CUBoulder.pdf}
	\end{minipage}
	\begin{minipage}{0.62\textwidth}
	\scriptsize
	\noindent Department of Chemical and Biological Engineering \hfill t 303-735-7860~~~~~~~~~~~~~~~~~~ \\
	\noindent 596 UCB \hfill f 303-492-8425~~~~~~~~~~~~~~~~~~ \\
	\noindent Boulder, Colorado 80309 \hfill michael.shirts@colorado.edu \\
	\end{minipage}
	\end{figure}
	
	%\noindent May 3\textsuperscript{rd}, 2019\\
	\noindent \today

	\noindent Dear Editors,\\
	
	\newcommand{\ManuscriptTitle}{Statistical Inference of Transport Mechanisms and
	Long Time Scale Behavior from Time Series of Solute Trajectories in 
	Nanostructured Membranes}
	
	We are submitting revisions to our manuscript ``\ManuscriptTitle'' by Benjamin 
	J. Coscia, Christopher P. Calderon and Michael R. Shirts, for consideration for 
	publication as a Research Article in the ``Machine Learning in Physical Chemistry"
	special issue of the Journal of Physical Chemistry, specifically in part B, section
	B4: Fluid Interfaces, Colloids, Polymers, Soft Matter, Surfactants, and Glassy Materials.

	We have received generally positive and constructive feedback
        from two reviewers which we have used to further strengthen
        our work. We have made many of the changes and provided
        rationales for not incorporating the other changes, as
        detailed in the attached ``Response to Reviewers''. We hope
        the mauscript is now acceptable for publication. For further
        description of the goals and scope of the research and suggested
        reviewers, see the cover letter included with the the previous version.

	
	%BJC: should I summarize the changes that were made here?

	% BJC: not sure if I should summarize again (this is from the original cover letter
	%In this article, we apply the sticky hierarchical Dirichlet process autoregressive
	%hidden Markov model (HDP-AR-HMM), a non-parametric Bayesian time series
	%classification technique, to solute center-of-mass trajectories generated from
	%5 $\mu s$ molecular dynamics (MD) simulations in a cross-linked inverted hexagonal
	%phase lyotropic liquid crystal (LLC) membrane in order to automatically detect a
	%variety of solute dynamical modes. We can use this information in order to 
	%better understand the complex transport mechanisms exhibited by the solutes in 
	%our simulation which will greatly increase the efficiency with which we 
	%design LLC membranes for solute-specific separations.
	
	%Our work represents a considerable advancement over our most recent study 
	%characterizing the same set of solute center-of-mass time series trajectories
	%(Coscia and Shirts, arXiv:2004.07905 [cond-mat.soft]). In that work, we used
	%our previously gained qualitative understanding of solute transport (Coscia 
	%and Shirts \textit{J. Phys. Chem. B}, 123, 6314--6330 (2019)) in order to 
	%formulate stochastic models which could be used to project solute behavior
	%on macroscopic time scales, thus allowing prediction of experimentally relevant
	%properties like membrane selectivity. Using the HDP-AR-HMM, we can analyze solute
	%trajectories with no prior knowledge of their dynamic behavior and easily produce
	%models which quantitatively describe their long time scale dynamics. Since the 
	%approach is non-parametric, the complexity of our data drives the complexity of
	%the model. We cluster the parameter sets in order to reduce the total space of 
	%dynamical states to a small set of distinct parameters which allow us to clearly
	%understand physical	mechanisms in terms of solute-membrane interactions. To our 
	%knowledge, this is the first application of the HDP-AR-HMM to MD data and it is
	%the first time this framework has been used to gain both mechanistic and kinetic
	%insight. 
	
	%BJC: I assume I can get rid of this?
	%\noindent Some suggestions for reviewers are:
	%\begin{enumerate}
	
	%	\item Francisco Hung has used molecular simulations to study nematic liquid crystals
%		and transport under confinement within various nanostructured materials 
%		(Northeastern University, 617-373-8619,\\ \href{mailto:f.hung@northeastern.edu}{f.hung@northeastern.edu}).
		
%		\item Mahesh Mahanthappa has experience modeling soft materials using molecular dynamics
%		including work on ion transport in various ordered liquid crystalline phases. 
%		(University of Minnesota, 612-625-4599, \href{mailto:maheshkm@umn.edu}{maheshkm@umn.edu}).
		
%		\item Stephen Paddison uses various levels of molecular simulation theory to understand
%		proton transport in nanostructured and charged polymer membranes. (The University of Tennessee 
%		Knoxville, 865-974-2026, \href{mailto:spaddison@utk.edu}{spaddison@utk.edu}).
		
%		\item Venkat Ganesan uses multiscale modeling to study solute and ion transport in complex
%		nanostructured morphologies in order to learn how to design novel materials. (The University
%		of Texas at Austin, 512-471-4856, \href{mailto:venkat@che.utexas.edu}{venkat@che.utexas.edu}).
		
		%BJC: from structure paper
%		\item Eric Jankowski has experience modeling soft materials using molecular dynamics.
%		(Boise State University, 208-426-5681, \href{mailto:ericjankowski@boisestate.edu}{ericjankowski@boisestate.edu})
		
%		\item Coray Colina uses various simulation techniques to understand structure-property relations
%		in functional materials including nanoporous materials for separation applications. (University
%		of Florida, 352-294-3488, \href{mailto:colina@chem.ufl.edu}{colina@chem.ufl.edu}).
		
%		\item Martin Lindén has use variational Bayesian treatments of hidden Markov models in
%		order to learn from thousands of single-particle trajectories.
%		
%		\item Steve Presse
%		
%	\end{enumerate}
	
	Please send correspondence regarding this paper to Michael R. Shirts (contact
	details in letterhead).\\	
	
	\noindent Sincerely,
	
	\noindent Michael R. Shirts \\
	\noindent Benjamin J. Coscia \\
	\noindent Christopher P. Calderon \\
	
\end{document}

% LocalWords:  BJC
